\documentclass[11pt,a4paper]{article}
\usepackage[utf8]{inputenc}
\usepackage[T1]{fontenc}
\usepackage[brazil]{babel}
\usepackage{geometry}
\usepackage{enumitem}
\usepackage{xcolor}
\usepackage{titlesec}
\usepackage{fancyhdr}
\usepackage{listings}
\usepackage{tcolorbox}
\usepackage{longtable}

\geometry{margin=2.5cm}

% Cores
\definecolor{headerblue}{RGB}{0,102,204}
\definecolor{sectiongray}{RGB}{100,100,100}
\definecolor{codebg}{RGB}{245,245,245}
\definecolor{loggreen}{RGB}{46,125,50}
\definecolor{logred}{RGB}{198,40,40}

% Estilo de seções
\titleformat{\section}{\Large\bfseries\color{headerblue}}{}{0em}{}
\titleformat{\subsection}{\large\bfseries\color{sectiongray}}{}{0em}{}

% Cabeçalho e rodapé
\pagestyle{fancy}
\fancyhf{}
\fancyhead[L]{\textbf{SIVIRA - Sistema de Logs}}
\fancyhead[R]{\thepage}
\renewcommand{\headrulewidth}{0.5pt}

\begin{document}

% Capa
\begin{titlepage}
    \centering
    \vspace*{3cm}

    {\Huge\bfseries Sistema Inteligente para Visualização\par}
    {\Huge\bfseries e Replanejamento de Atividades\par}
    \vspace{0.5cm}
    {\Large\bfseries (SIVIRA)\par}

    \vspace{2cm}
    {\LARGE\bfseries Documentação do Sistema de Logs\par}

    \vspace{2cm}
    {\large Módulo: \texttt{logs/}\par}

    \vspace{1cm}
    {\large Rastreamento Completo de Execução\par}

    \vfill

    {\large \today\par}
\end{titlepage}

\tableofcontents
\newpage

% Introdução
\section{Introdução}

O sistema de logs do SIVIRA é uma infraestrutura completa e hierárquica de rastreamento de execução que registra todos os aspectos da operação do sistema de produção. Os logs são organizados em 10 categorias principais, cada uma com propósito específico e formato padronizado.

\subsection{Objetivos do Sistema de Logs}

\begin{itemize}
    \item \textbf{Rastreabilidade}: Registro completo de todas as operações executadas
    \item \textbf{Auditoria}: Histórico detalhado para análise posterior
    \item \textbf{Depuração}: Identificação rápida de problemas e erros
    \item \textbf{Otimização}: Análise de desempenho e gargalos
    \item \textbf{Validação}: Verificação de conformidade com requisitos
\end{itemize}

\subsection{Estrutura de Diretórios}

\begin{tcolorbox}[colback=codebg,colframe=headerblue,title=Hierarquia de Pastas de Logs]
\begin{verbatim}
logs/
├── agendas/                 📅 Agendas de equipamentos
├── equipamentos/            🔧 Alocações de equipamentos
├── equipamentos_detalhados/ 📊 Ocupações detalhadas
├── equipamentos_reservados/ 🔒 Reservas de equipamentos
├── erros/                   ❌ Erros e falhas
├── execucoes/              🚀 Histórico de execuções
├── funcionarios/           👷 Alocações de funcionários
├── restricoes/             ⚖️ Restrições aplicadas
├── temp/                   📝 Arquivos temporários
└── tipos_funcionarios_requeridos/ 👥 Requisitos profissionais
\end{verbatim}
\end{tcolorbox}

\newpage

\section{Pastas de Logs}

\subsection{logs/agendas/ - Agendas de Equipamentos}

\subsubsection{Propósito}

Armazena visualizações completas das agendas de equipamentos após execução de pedidos, incluindo todas as ocupações, configurações e detalhes técnicos.

\subsubsection{Formato de Arquivo}

\texttt{agenda\_equipamentos\_completa\_YYYYMMDD\_HHMMSS.log}

\subsubsection{Conteúdo}

\begin{itemize}[leftmargin=2cm]
    \item \textbf{Cabeçalho}: Data/hora de geração, identificação do sistema
    \item \textbf{Por equipamento}:
    \begin{itemize}
        \item Nome e tipo do equipamento
        \item Configurações específicas (temperatura, velocidade, etc.)
        \item Todas as ocupações com horários
        \item Detalhes técnicos (bocas de fogão, níveis de forno, etc.)
    \end{itemize}
    \item \textbf{Estatísticas}: Número total de equipamentos, taxas de utilização
\end{itemize}

\subsubsection{Módulo Gerador}

\begin{itemize}
    \item \texttt{examples/producao/*\_com\_agenda.py}
    \item Método: \texttt{\_salvar\_agenda\_completa\_em\_arquivo()}
\end{itemize}

\subsubsection{Quando é Criado}

Ao final da execução de pedidos em scripts que incluem funcionalidade de agenda (coxinhas, folhados, múltiplos produtos).

\subsection{logs/equipamentos/ - Alocações de Equipamentos}

\subsubsection{Propósito}

Registra todas as alocações de equipamentos para atividades de produção, permitindo rastrear qual equipamento foi usado em cada atividade.

\subsubsection{Formato de Arquivo}

\texttt{ordem: [ID\_ORDEM] | pedido: [ID\_PEDIDO].log}

\subsubsection{Conteúdo}

Formato de linha (separado por \texttt{|}):
\begin{verbatim}
ID_ORDEM | ID_PEDIDO | ID_ATIVIDADE | NOME_ITEM |
NOME_ATIVIDADE | NOME_EQUIPAMENTO | HORA_INICIO | HORA_FIM
\end{verbatim}

\textbf{Exemplo}:
\begin{verbatim}
1 | 1 | 10011 | Pão Francês | pesagem_de_massas |
Balança Digital 1 | 02:12 | 02:24
\end{verbatim}

\subsubsection{Módulos Geradores}

\begin{itemize}
    \item \texttt{utils/logs/logger\_de\_atividades.py}
    \item \texttt{utils/logs/gerenciador\_logs.py}
    \item Função: \texttt{registrar\_log\_equipamentos()}
\end{itemize}

\subsubsection{Quando é Criado}

Durante a execução de cada atividade que requer equipamento, em modo append para construir histórico completo do pedido.

\subsection{logs/equipamentos\_detalhados/ - Ocupações Detalhadas}

\subsubsection{Propósito}

Fornece visão extremamente detalhada das ocupações de equipamentos, incluindo informações granulares como bocas específicas de fogão, níveis de forno, frações de bancada, etc.

\subsubsection{Formato de Arquivo}

\texttt{ocupacoes\_detalhadas\_ordem\_[ID\_ORDEM]\_pedidos\_[LISTA\_PEDIDOS]\_YYYYMMDD\_HHMMSS.log}

\subsubsection{Conteúdo}

Para cada equipamento:
\begin{itemize}[leftmargin=2cm]
    \item \textbf{Identificação}: Nome, tipo, ID
    \item \textbf{Ocupações simples}: Lista cronológica de usos
    \item \textbf{Ocupações por nível/boca/fração}:
    \begin{itemize}
        \item Fornos: ocupações por nível (1-5)
        \item Fogões: ocupações por boca (1-6)
        \item Bancadas: ocupações por fração
        \item Fritadeiras: ocupações por fração
    \end{itemize}
    \item \textbf{Detalhes técnicos}: Configurações, capacidades, restrições
\end{itemize}

\subsubsection{Módulos Geradores}

\begin{itemize}
    \item \texttt{utils/logs/logger\_ocupacao\_detalhada.py}
    \item \texttt{utils/logs/capturador\_ocupacoes\_equipamentos.py}
    \item Classes: \texttt{LoggerOcupacaoDetalhada}, \texttt{CapturadorOcupacoesEquipamentos}
\end{itemize}

\subsubsection{Quando é Criado}

Ao final da execução de pedidos, capturando estado completo de todos os equipamentos utilizados.

\subsection{logs/equipamentos\_reservados/ - Reservas de Equipamentos}

\subsubsection{Propósito}

Registra reservas temporárias de equipamentos durante o processo de alocação, útil para debug de algoritmos de scheduling.

\subsubsection{Formato de Arquivo}

\texttt{*.log} (formato variado conforme necessidade)

\subsubsection{Conteúdo}

Informações sobre:
\begin{itemize}
    \item Equipamentos marcados como reservados
    \item Período de reserva
    \item Atividade que solicitou a reserva
    \item Status da reserva (confirmada, liberada, expirada)
\end{itemize}

\subsubsection{Quando é Criado}

Durante processos de otimização e alocação que utilizam sistema de reservas.

\subsection{logs/erros/ - Erros e Falhas}

\subsubsection{Propósito}

Sistema centralizado e estruturado para registro de todos os erros, falhas e exceções do sistema, com contexto completo para análise.

\subsubsection{Formato de Arquivo}

\begin{itemize}
    \item \textbf{JSON estruturado}: \texttt{[TIPO\_ERRO]\_[ORDEM]\_[PEDIDO]\_[ATIVIDADE]\_[TIMESTAMP].json}
    \item \textbf{Log de compatibilidade}: \texttt{ordem: [ID] | pedido: [ID].log}
\end{itemize}

\subsubsection{Conteúdo do JSON}

\begin{tcolorbox}[colback=codebg,colframe=logred,title=Estrutura de Erro JSON]
\begin{lstlisting}[language=json]
{
  "timestamp": "2025-10-06T18:00:00",
  "identificacao": {
    "id_ordem": 1,
    "id_pedido": 1,
    "id_atividade": 10011,
    "nome_atividade": "pesagem_de_massas"
  },
  "erro": {
    "tipo": "TEMPORAL_ALLOCATION_ERROR",
    "detalhes": {...},
    "nivel_impacto": "ALTO"
  },
  "contexto": {...},
  "sistema": {
    "versao_log": "2.0",
    "gerado_por": "ErrorLogger"
  }
}
\end{lstlisting}
\end{tcolorbox}

\subsubsection{Tipos de Erro Registrados}

\begin{itemize}[leftmargin=2cm]
    \item \textbf{TEMPORAL\_ALLOCATION\_ERROR}: Falhas em alocação temporal
    \item \textbf{DEADLINE\_ALLOCATION\_ERROR}: Violações de prazo
    \item \textbf{QUANTITY\_ALLOCATION\_ERROR}: Problemas de quantidade
    \item \textbf{CONFIGURATION\_ERROR}: Erros de configuração
    \item \textbf{TIMING\_ERROR}: Problemas de timing
    \item \textbf{RESOURCE\_ERROR}: Falta de recursos
\end{itemize}

\subsubsection{Níveis de Impacto}

\begin{itemize}
    \item \textbf{CRÍTICO}: Sistema não pode continuar
    \item \textbf{ALTO}: Pedido/atividade falharam
    \item \textbf{MÉDIO}: Degradação de performance
    \item \textbf{BAIXO}: Avisos informativos
\end{itemize}

\subsubsection{Módulos Geradores}

\begin{itemize}
    \item \texttt{utils/logs/error\_logger\_utils.py}
    \item \texttt{utils/logs/temporal\_allocation\_logger.py}
    \item \texttt{utils/logs/deadline\_allocation\_logger.py}
    \item \texttt{utils/logs/timing\_exceptions.py}
    \item \texttt{utils/logs/quantity\_exceptions.py}
    \item Classe principal: \texttt{ErrorLogger}
\end{itemize}

\subsection{logs/execucoes/ - Histórico de Execuções}

\subsubsection{Propósito}

Mantém histórico de alto nível de todas as execuções de pedidos, servindo como índice para navegação nos demais logs.

\subsubsection{Formato de Arquivo}

\texttt{execucao\_[TIPO]\_[TIMESTAMP].log}

\textbf{Exemplos}:
\begin{itemize}
    \item \texttt{execucao\_pedidos\_20251006\_180000.log}
    \item \texttt{execucao\_coxinhas\_20251006\_180000.log}
\end{itemize}

\subsubsection{Conteúdo}

\begin{itemize}[leftmargin=2cm]
    \item \textbf{Cabeçalho}: Data/hora, versão do sistema
    \item \textbf{Configuração}: Parâmetros de execução
    \item \textbf{Pedidos processados}: Lista com status
    \item \textbf{Estatísticas}:
    \begin{itemize}
        \item Tempo total de execução
        \item Número de atividades processadas
        \item Taxa de sucesso
        \item Recursos utilizados
    \end{itemize}
    \item \textbf{Referências}: Links para outros logs relacionados
\end{itemize}

\subsubsection{Módulos Geradores}

\begin{itemize}
    \item \texttt{utils/logs/gerenciador\_logs.py}
    \item \texttt{utils/logs/gerenciador\_logs\_fases.py}
\end{itemize}

\subsection{logs/funcionarios/ - Alocações de Funcionários}

\subsubsection{Propósito}

Registra todas as alocações de funcionários para atividades de produção, rastreando quem fez o quê e quando.

\subsubsection{Formato de Arquivo}

\texttt{ordem: [ID\_ORDEM] | pedido: [ID\_PEDIDO].log}

\subsubsection{Conteúdo}

Formato de linha (separado por \texttt{|}):
\begin{verbatim}
ID_ORDEM | ID_PEDIDO | ID_ATIVIDADE | NOME_ITEM |
NOME_ATIVIDADE | NOME_FUNCIONARIO [STATUS] |
HORA_INICIO [DATA] | HORA_FIM [DATA]
\end{verbatim}

\textbf{Status possíveis}:
\begin{itemize}
    \item \texttt{✅} - Funcionário alocado com sucesso
    \item \texttt{❌} - Funcionário indisponível (com tipos necessários)
\end{itemize}

\textbf{Exemplo}:
\begin{verbatim}
1 | 1 | 10011 | Pão Francês | pesagem_de_massas |
Funcionário 1 ✅ | 02:12 [11/10/2025] | 02:24 [11/10/2025]
\end{verbatim}

\subsubsection{Módulos Geradores}

\begin{itemize}
    \item \texttt{utils/logs/gerenciador\_logs.py}
    \item \texttt{services/gestores/funcionarios/gestor\_funcionarios.py}
    \item Função: \texttt{registrar\_log\_funcionarios()}
\end{itemize}

\subsubsection{Quando é Criado}

Durante a alocação automática de funcionários para atividades de pedidos, registrando tanto sucessos quanto falhas.

\subsection{logs/restricoes/ - Restrições Aplicadas}

\subsubsection{Propósito}

Documenta todas as restrições aplicadas durante o processo de scheduling, permitindo entender decisões de alocação.

\subsubsection{Formato de Arquivo}

\texttt{ordem\_[ID\_ORDEM]\_restricoes.json}

\subsubsection{Conteúdo do JSON}

\begin{tcolorbox}[colback=codebg,colframe=headerblue,title=Estrutura de Restrições]
\begin{lstlisting}[language=json]
{
  "ordem": 1,
  "pedidos": [1, 2, 3],
  "restricoes": [
    {
      "id_pedido": 1,
      "id_atividade": 10011,
      "tipo_restricao": "TEMPORAL",
      "detalhes": {
        "inicio_minimo": "2025-10-11T02:00:00",
        "inicio_maximo": "2025-10-11T03:00:00"
      },
      "aplicada": true,
      "impacto": "Limitou janela de execucao"
    }
  ],
  "timestamp": "2025-10-06T18:00:00"
}
\end{lstlisting}
\end{tcolorbox}

\subsubsection{Tipos de Restrição}

\begin{itemize}[leftmargin=2cm]
    \item \textbf{TEMPORAL}: Restrições de horário/prazo
    \item \textbf{PRECEDENCIA}: Ordem de execução entre atividades
    \item \textbf{RECURSO}: Disponibilidade de equipamentos/funcionários
    \item \textbf{CAPACIDADE}: Limites de produção
    \item \textbf{QUALIDADE}: Requisitos de processo
\end{itemize}

\subsubsection{Módulos Geradores}

\begin{itemize}
    \item \texttt{utils/logs/registrador\_restricoes.py}
    \item Classe: \texttt{RegistradorRestricoes}
\end{itemize}

\subsection{logs/temp/ - Arquivos Temporários}

\subsubsection{Propósito}

Armazena logs temporários durante execução, úteis para debug em tempo real mas que podem ser descartados após processamento.

\subsubsection{Conteúdo}

\begin{itemize}
    \item Logs intermediários de processos longos
    \item Dumps de estado para recuperação
    \item Snapshots de dados em processamento
    \item Caches temporários
\end{itemize}

\subsubsection{Ciclo de Vida}

Arquivos nesta pasta podem ser automaticamente limpos:
\begin{itemize}
    \item No início de nova execução
    \item Após conclusão bem-sucedida
    \item Por rotinas de manutenção
\end{itemize}

\subsection{logs/tipos\_funcionarios\_requeridos/ - Requisitos Profissionais}

\subsubsection{Propósito}

Documenta os tipos profissionais e funcionários necessários para cada atividade de cada pedido, servindo como base para alocação de mão de obra.

\subsubsection{Formato de Arquivo}

\texttt{ordem: [ID\_ORDEM] | pedido: [ID\_PEDIDO].log}

\subsubsection{Conteúdo}

Para cada atividade do pedido:

\begin{tcolorbox}[colback=codebg,colframe=loggreen,title=Formato de Requisito]
\begin{verbatim}
🔸 ATIVIDADE 10011: pesagem_de_massas_para_pao_frances
⏰ Horário: 02:12 [11/10/2025] - 02:24 [11/10/2025]
👥 Funcionários necessários: 2
👨‍💼 Tipos permitidos: PADEIRO, AUXILIAR_DE_PADEIRO
🎯 Prioridades (FIPs): PADEIRO: 2, AUXILIAR_DE_PADEIRO: 1
────────────────────────────────────────────────────
\end{verbatim}
\end{tcolorbox}

\textbf{Informações detalhadas}:
\begin{itemize}[leftmargin=2cm]
    \item \textbf{ID e nome da atividade}
    \item \textbf{Janela temporal}: Horário de início e fim com data completa
    \item \textbf{Quantidade}: Número de funcionários necessários
    \item \textbf{Tipos profissionais}: Lista de profissões aceitas
    \item \textbf{FIP}: Fator de Importância do Profissional para cada tipo
\end{itemize}

\subsubsection{Módulos Geradores}

\begin{itemize}
    \item \texttt{utils/logs/registrador\_funcionarios.py}
    \item Classe: \texttt{RegistradorFuncionarios}
    \item Chamado por: \texttt{models/atividades/atividade\_modular.py}
\end{itemize}

\subsubsection{Quando é Criado}

Durante a execução de atividades que requerem mão de obra, antes da alocação efetiva de funcionários.

\subsubsection{Uso}

Este log serve como entrada para o \texttt{GestorFuncionarios} que:
\begin{enumerate}
    \item Lê os requisitos deste arquivo
    \item Identifica funcionários disponíveis e compatíveis
    \item Aplica algoritmo de priorização baseado em FIP
    \item Aloca funcionários e registra em \texttt{logs/funcionarios/}
\end{enumerate}

\newpage

\section{Módulos Responsáveis}

\subsection{utils/logs/gerenciador\_logs.py}

Módulo principal de gerenciamento de logs.

\textbf{Responsabilidades}:
\begin{itemize}
    \item Limpeza automática de logs na inicialização
    \item Funções para registro em equipamentos e funcionários
    \item Gerenciamento de pastas e arquivos
    \item Utilitários de formatação
\end{itemize}

\textbf{Funções principais}:
\begin{itemize}[leftmargin=2cm]
    \item \texttt{limpar\_logs\_inicializacao()}: Limpa pastas de log
    \item \texttt{registrar\_log\_equipamentos()}: Registra uso de equipamento
    \item \texttt{registrar\_log\_funcionarios()}: Registra alocação de funcionário
    \item \texttt{apagar\_logs\_por\_pedido\_e\_ordem()}: Remove logs específicos
\end{itemize}

\subsection{utils/logs/error\_logger\_utils.py}

Sistema centralizado de logging de erros.

\textbf{Classe principal}: \texttt{ErrorLogger}

\textbf{Características}:
\begin{itemize}
    \item Formato JSON estruturado
    \item Níveis de impacto (CRÍTICO, ALTO, MÉDIO, BAIXO)
    \item Contexto completo do erro
    \item Compatibilidade com sistema legado
\end{itemize}

\subsection{utils/logs/logger\_ocupacao\_detalhada.py}

Gerador de logs extremamente detalhados de ocupações.

\textbf{Classe principal}: \texttt{LoggerOcupacaoDetalhada}

\textbf{Suporta}:
\begin{itemize}
    \item Equipamentos com níveis (fornos)
    \item Equipamentos com bocas (fogões)
    \item Equipamentos com frações (bancadas, fritadeiras)
    \item Equipamentos simples
\end{itemize}

\subsection{utils/logs/registrador\_funcionarios.py}

Registra requisitos de funcionários para atividades.

\textbf{Classe principal}: \texttt{RegistradorFuncionarios}

\textbf{Características}:
\begin{itemize}
    \item Extrai tipos profissionais necessários
    \item Calcula FIP para priorização
    \item Formata requisitos de forma legível
    \item Integra com sistema de alocação
\end{itemize}

\subsection{utils/logs/registrador\_restricoes.py}

Documenta restrições aplicadas durante scheduling.

\textbf{Classe principal}: \texttt{RegistradorRestricoes}

\textbf{Tipos de restrição}:
\begin{itemize}
    \item Temporais
    \item De precedência
    \item De recurso
    \item De capacidade
\end{itemize}

\newpage

\section{Fluxo de Logging}

\subsection{Durante Execução de Pedido}

\begin{enumerate}
    \item \textbf{Início}: \texttt{logs/execucoes/} recebe cabeçalho
    \item \textbf{Por atividade}:
    \begin{itemize}
        \item Restrições → \texttt{logs/restricoes/}
        \item Requisitos funcionários → \texttt{logs/tipos\_funcionarios\_requeridos/}
        \item Alocação equipamentos → \texttt{logs/equipamentos/}
    \end{itemize}
    \item \textbf{Alocação de RH}:
    \begin{itemize}
        \item GestorFuncionarios lê \texttt{logs/tipos\_funcionarios\_requeridos/}
        \item Aloca funcionários
        \item Registra em \texttt{logs/funcionarios/}
    \end{itemize}
    \item \textbf{Erros}: Qualquer falha → \texttt{logs/erros/}
    \item \textbf{Finalização}:
    \begin{itemize}
        \item Ocupações detalhadas → \texttt{logs/equipamentos\_detalhados/}
        \item Agenda completa → \texttt{logs/agendas/}
        \item Estatísticas → \texttt{logs/execucoes/}
    \end{itemize}
\end{enumerate}

\subsection{Limpeza Automática}

No início de cada execução do menu principal:

\begin{enumerate}
    \item Limpa \texttt{logs/funcionarios/}
    \item Limpa \texttt{logs/equipamentos/}
    \item Limpa \texttt{logs/erros/}
    \item Limpa \texttt{logs/execucoes/}
    \item Limpa \texttt{logs/restricoes/}
    \item Limpa \texttt{logs/equipamentos\_detalhados/}
    \item Remove \texttt{data/pedidos/pedidos\_salvos.json}
\end{enumerate}

\textbf{Pastas preservadas}:
\begin{itemize}
    \item \texttt{logs/agendas/} - Histórico valioso
    \item \texttt{logs/tipos\_funcionarios\_requeridos/} - Gerado sob demanda
\end{itemize}

\newpage

\section{Formatos e Convenções}

\subsection{Timestamps}

\textbf{Em nomes de arquivo}:
\begin{verbatim}
YYYYMMDD_HHMMSS
Exemplo: 20251006_180000
\end{verbatim}

\textbf{Em conteúdo (JSON)}:
\begin{verbatim}
ISO 8601: YYYY-MM-DDTHH:MM:SS
Exemplo: 2025-10-06T18:00:00
\end{verbatim}

\textbf{Em logs texto}:
\begin{verbatim}
HH:MM [DD/MM/YYYY]
Exemplo: 18:00 [06/10/2025]
\end{verbatim}

\subsection{Separadores}

\textbf{Logs estruturados}:
\begin{itemize}
    \item Campo: \texttt{|} (pipe)
    \item Espaçamento: \texttt{ | } (pipe com espaços)
\end{itemize}

\textbf{Nomes de arquivo}:
\begin{itemize}
    \item Palavras: \texttt{\_} (underscore)
    \item Componentes: \texttt{ordem\_1\_pedido\_1}
\end{itemize}

\subsection{Encoding}

Todos os arquivos utilizam:
\begin{itemize}
    \item \textbf{Encoding}: UTF-8
    \item \textbf{Line ending}: Unix (LF)
    \item \textbf{BOM}: Não utilizado
\end{itemize}

\newpage

\section{Estatísticas}

\subsection{Resumo Geral}

\begin{tabular}{lr}
\textbf{Métrica} & \textbf{Valor} \\
\hline
Total de pastas de logs & 10 \\
Módulos geradores de logs & 17 \\
Formatos de arquivo & 2 (LOG, JSON) \\
Tipos de erro categorizados & 6+ \\
Níveis de impacto & 4 \\
\hline
\end{tabular}

\subsection{Pastas por Categoria}

\begin{tabular}{lll}
\textbf{Categoria} & \textbf{Pastas} & \textbf{Propósito} \\
\hline
Recursos & 2 & equipamentos, funcionarios \\
Detalhamento & 3 & equipamentos\_detalhados, agendas, tipos\_funcionarios\_requeridos \\
Controle & 2 & restricoes, execucoes \\
Erros & 1 & erros \\
Temporário & 2 & temp, equipamentos\_reservados \\
\hline
\textbf{Total} & \textbf{10} & \\
\end{tabular}

\subsection{Módulos de Logging}

\begin{longtable}{lp{8cm}}
\textbf{Módulo} & \textbf{Responsabilidade} \\
\hline
\endfirsthead
\textbf{Módulo} & \textbf{Responsabilidade} \\
\hline
\endhead

gerenciador\_logs.py & Coordenação geral, limpeza, equipamentos, funcionários \\
error\_logger\_utils.py & Sistema centralizado de erros \\
logger\_ocupacao\_detalhada.py & Ocupações detalhadas de equipamentos \\
capturador\_ocupacoes\_equipamentos.py & Captura estado de equipamentos \\
registrador\_funcionarios.py & Requisitos profissionais \\
registrador\_restricoes.py & Restrições de scheduling \\
temporal\_allocation\_logger.py & Erros temporais \\
deadline\_allocation\_logger.py & Erros de prazo \\
timing\_exceptions.py & Exceções de timing \\
quantity\_exceptions.py & Exceções de quantidade \\
logger\_de\_atividades.py & Logs de atividades \\
gerenciador\_logs\_fases.py & Logs por fase \\
log\_subprodutos\_agrupados.py & Agrupamento de subprodutos \\
formatador\_logs\_limpo.py & Formatação de logs \\
formatador\_timing\_limpo.py & Formatação de timing \\
timing\_logger.py & Logs de timing \\
quantity\_logger.py & Logs de quantidade \\
\hline
\end{longtable}

\newpage

\section{Boas Práticas}

\subsection{Para Desenvolvedores}

\begin{enumerate}
    \item \textbf{Sempre use encoding UTF-8} ao abrir arquivos de log
    \item \textbf{Adicione timestamps} em todas as entradas de log
    \item \textbf{Use ErrorLogger} para erros estruturados
    \item \textbf{Documente formato} de novos tipos de log
    \item \textbf{Mantenha compatibilidade} com ferramentas de análise
\end{enumerate}

\subsection{Para Análise}

\begin{enumerate}
    \item \textbf{Comece por logs/execucoes/} para visão geral
    \item \textbf{Use logs/erros/} para investigar problemas
    \item \textbf{Cruze dados} entre equipamentos e funcionários
    \item \textbf{Valide restrições} com logs/restricoes/
    \item \textbf{Analise ocupações} com logs/equipamentos\_detalhados/
\end{enumerate}

\subsection{Manutenção}

\begin{enumerate}
    \item \textbf{Limpe logs antigos} periodicamente
    \item \textbf{Arquive históricos} importantes
    \item \textbf{Monitore tamanho} das pastas de log
    \item \textbf{Verifique permissões} de escrita
    \item \textbf{Faça backup} antes de limpezas massivas
\end{enumerate}

\newpage

\section{Conclusão}

O sistema de logs do SIVIRA fornece rastreabilidade completa e estruturada de todas as operações do sistema de produção. Com 10 categorias de logs, 17 módulos geradores e suporte a múltiplos formatos, o sistema permite:

\begin{itemize}
    \item \textbf{Auditoria completa}: Cada operação é rastreável
    \item \textbf{Debug eficiente}: Erros contextualizados e categorizados
    \item \textbf{Análise de desempenho}: Logs detalhados de ocupações
    \item \textbf{Otimização}: Identificação de gargalos e oportunidades
    \item \textbf{Conformidade}: Histórico para validação de processos
\end{itemize}

\subsection{Próximos Passos}

Possíveis melhorias futuras:

\begin{enumerate}
    \item Visualização gráfica de logs
    \item Sistema de alertas baseado em logs
    \item Análise automática de padrões
    \item Compressão de logs antigos
    \item Dashboard de monitoramento em tempo real
    \item Exportação para formatos analíticos (CSV, Parquet)
\end{enumerate}

\end{document}
